\section{Wie benutze ich das Glossar?}
\label{sec:glossar}

Mit \lstinline|\printglossaries| wird das Glossar eingefügt. Glossar-Einträge können in \textit{glossary.tex} gepflegt werden.

\begin{lstlisting}[language=TeX, caption=Beispiel für zwei Glossareinträge]
\newglossaryentry{latex}
{
        name=\LaTeX,
        description={Eine Markup-Sprache, geeignet für wissenschaftliche Arbeiten}
}

\newacronym{ctan}{CTAN}{Comprehensive \TeX Archive Network}
\end{lstlisting}

Nicht alle Glossareinträge erscheinen sofort in den Glossar- oder Abkürzungslisten, denn sie müssen erst verwendet werden.

Akronyme wie \gls{ctan} können mit \lstinline|\gls{ctan}| ausgegeben werden. Das Gleiche gilt auch für Glossareinträge wie \gls{latex}.

Weitere Informationen unter \href{https://ctan.org/pkg/glossaries}{CTAN Glossaries}.